% !TEX root = ../main.tex
\subsection*{Example --- Equality Lebesgue-Almost Everywhere}%
\manualtoc{subsection}{\texorpdfstring{Example --- Equality Lebesgue-Almost Everywhere}{Example --- Equality Lebesgue-Almost Everywhere}}{studiabile-nav-1004}%

Let $S=[0,1]$, $\Sigma=\mathcal B([0,1])$, $\mu=\mathrm{Leb}$, and consider
\[
f=\mathbf1_{[\frac12,1]},\qquad
 g=\mathbf1_{(\frac12,1]}.
\]
The functions differ only at $s=\frac12$, so
\[
\{s\in S:f(s)\ne g(s)\}=\bigl\{\tfrac12\bigr\},
\qquad \mathrm{Leb}\bigl(\{\tfrac12\}\bigr)=0.
\]
Hence $f=g$ $\mathrm{Leb}$-a.e. and, by HW~44(i),
\[
\mu_0(f)
=\mathrm{Leb}\bigl([\tfrac12,1]\bigr)
=\tfrac12
=\mathrm{Leb}\bigl((\tfrac12,1]\bigr)
=\mu_0(g).
\]
The equalities follow from the definition of the integral of an indicator function and the length of the two intervals.

\subsection*{Definition --- Integral of a Positive Measurable Function}%
\manualtoc{subsection}{\texorpdfstring{Definition --- Integral of a Positive Measurable Function}{Definition --- Integral of a Positive Measurable Function}}{studiabile-nav-1005}%

In this chapter, $(m\Sigma)^+$ includes all measurable maps $f:S\to[0,+\infty]$, where measurability is understood through the real threshold tests. Simple functions in $SF^+$ still have only finitely many finite values.

For signed functions, $m\Sigma$ retains the real-valued meaning of Definition~3.1 unless stated otherwise.

Let $(S,\Sigma,\mu)$ be a measure space. For any $f\in(m\Sigma)^+$, define
\[
\mu(f):=\int_S f\,d\mu
:=\sup\bigl\{\mu_0(h):h\in SF^+,\ h\le f\bigr\}.
\]
Thus the integral of $f$ is the supremum of the integrals of all nonnegative simple functions lying below $f$.

\subsection*{Observation --- Agreement with the Simple Integral}%
\manualtoc{subsection}{\texorpdfstring{Observation --- Agreement with the Simple Integral}{Observation --- Agreement with the Simple Integral}}{studiabile-nav-1006}%

If $f\in SF^+$, then
\[
\mu(f)=\mu_0(f).
\]
Indeed, $f$ itself belongs to the family over which the supremum is taken, so $\mu(f)\ge\mu_0(f)$. Conversely, for every $h\in SF^+$ with $h\le f$, monotonicity of $\mu_0$ (HW~44(iii)) gives $\mu_0(h)\le\mu_0(f)$. Taking the supremum yields $\mu(f)\le\mu_0(f)$, proving the equality.

\subsection*{HW 45 --- Monotonicity of the Positive Lebesgue Integral}%
\manualtoc{subsection}{\texorpdfstring{HW 45 --- Monotonicity of the Positive Lebesgue Integral}{HW 45 --- Monotonicity of the Positive Lebesgue Integral}}{studiabile-nav-1007}%

\studyheading{Statement.} Prove that if $f,g\in(m\Sigma)^+$ satisfy $f\le g$, then
\[
\mu(f)\le\mu(g).
\]
